% id: 138-24
% book: 베이직쎈 확률과 통계 · 07 통계적 추정
% section: 유형 089 모평균에 대한 신뢰구간의 성질
% figure: none
% answer: 
% answer_source: 
% variant_level: 0 (원본 전사)
% 이 파일은 build.mjs 가 items.json 에서 만든다. 고칠 때는 items.json 을 고친다.
\smash{\raisebox{1.5mm}{\dmkichul{베이직쎈 확률과 통계 138쪽 24번}}}
정규분포를 따르는 모집단에서 표본을 임의추출하여 모평균을 추정할 때, 옳은 것만을 \textbf{보기}에서 있는 대로 고른 것은? \bogi{ㄱ. 표본의 크기가 일정할 때, 신뢰도를 낮추면 신뢰구간의 길이는 짧아진다.\\ ㄴ. 신뢰도가 일정할 때, 표본의 크기를 9배 하면 신뢰구간의 길이는 3배가 된다.\\ ㄷ. 신뢰도를 낮추면서 표본의 크기를 크게 하면 신뢰구간의 길이는 짧아진다.}
\dmchoicesauto{}{ㄱ}{ㄴ}{ㄷ}{ㄱ, ㄴ}{ㄱ, ㄷ}
